\documentclass[11pt,a4paper]{article}
\usepackage[utf8]{inputenc}
\usepackage{amsmath,amssymb,amsthm}
\usepackage{geometry}
\geometry{margin=1in}
\usepackage{hyperref}

\newtheorem{theorem}{Theorem}[section]
\newtheorem{lemma}[theorem]{Lemma}
\newtheorem{definition}[theorem]{Definition}
\newtheorem{conjecture}[theorem]{Conjecture}

\title{Rigorous Deconstruction and Partial Resolution Strategy for the Erdős-Sós Conjecture}
\author{Charles EDOU NZE\thanks{Charles EDOU NZE, chercheur indépendant}}
\date{\today}

\begin{document}
\maketitle

\begin{abstract}
We present a rigorous axiomatic deconstruction and partial resolution strategy for the Erdős-Sós conjecture, which asserts that any graph with average degree strictly greater than $k-1$ must contain every tree on $k+1$ vertices as a subgraph. We introduce contextual literature bounds, formal structural definitions for lean autoformalization, and provide zero-ellipse proofs for key lemmas involving degree sequences and minimal counterexamples.
\end{abstract}

\tableofcontents

\section{Axiomatic Definitions and Problem Statement}
The Erdős-Sós conjecture (1962) remains one of the central problems in extremal graph theory. We begin by formalizing the definitions.

\begin{definition}[Graph and Tree]
Let $G = (V, E)$ be a simple, undirected graph where $V$ is a finite set of vertices and $E \subseteq \binom{V}{2}$ is the set of edges. A tree $T$ is a connected acyclic graph.
\end{definition}

\begin{definition}[Average Degree]
The average degree of a graph $G = (V, E)$ is given by $d(G) = \frac{2|E|}{|V|}$.
\end{definition}

\begin{conjecture}[Erdős-Sós Conjecture]
Let $k \in \mathbb{N}$ with $k \geq 1$. If $G$ is a graph with average degree $d(G) > k - 1$, then $G$ contains every tree $T$ on $k+1$ vertices as a subgraph.
\end{conjecture}

\section{Contextual Literature Research}
A literature review reveals analogies with related combinatorial problems. For instance, recent research on generalized connectivity and Turán-type bounds on bipartite graphs explore dense extremal structures. Constructing disjoint Steiner trees and investigating threshold graphs provide bounds on substructures avoiding specific components, similar in spirit to finding trees in graphs of bounded average degree. The classic Ajtai-Komlós-Szemerédi theorem provides related asymptotic bounds, and our approach aims to leverage structural decomposition analogous to those used in dense graph limit theories.

\section{Autoformalization Architecture}
The proof structure is designed to facilitate translation into proof assistants such as Lean 4. We specify the types and variables explicitly.

\begin{verbatim}
-- Lean 4 Architecture Specification
import Mathlib.Combinatorics.SimpleGraph.Basic
import Mathlib.Combinatorics.SimpleGraph.Connectivity

universe u
variable {V : Type u} [Fintype V] [DecidableEq V]

def AverageDegree (G : SimpleGraph V) : \mathbb{Q} :=
  (2 * G.edgeFinset.card) / Fintype.card V

def IsTree (G : SimpleGraph V) : Prop :=
  G.Connected \land G.IsAcyclic

theorem ErdosSos (k : \mathbb{N}) (G : SimpleGraph V) (T : SimpleGraph V)
  (hG : AverageDegree G > k - 1)
  (hT_tree : IsTree T)
  (hT_verts : Fintype.card V = k + 1) :
  \exists f : V \hookrightarrow V, \forall v w : V, T.Adj v w → G.Adj (f v) (f w) :=
  sorry
\end{verbatim}

\section{Lemmas and Zero-Ellipse Proofs}

We decompose the problem into several lemmas.

\begin{lemma}[Minimum Degree Reduction]
If there exists a counterexample $G$ to the Erdős-Sós conjecture for a given $k$, there exists a minimal counterexample $G^*$ such that the minimum degree $\delta(G^*) \geq \frac{k}{2}$.
\end{lemma}

\begin{proof}
Let $G = (V, E)$ be a graph with $d(G) > k-1$ such that $G$ does not contain a specific tree $T$ on $k+1$ vertices. Suppose for the sake of contradiction that there is no subgraph $H \subseteq G$ with $d(H) > k-1$ and minimum degree $\delta(H) \geq \frac{k}{2}$.

We construct a sequence of subgraphs $G_0, G_1, \dots$ recursively. Let $G_0 = G$.
For $i \geq 0$, if $G_i$ has a vertex $v_i$ with degree $d_{G_i}(v_i) \leq \frac{k-1}{2}$, we obtain $G_{i+1}$ by deleting $v_i$ and its incident edges.
At each step, the number of edges removed is at most $\frac{k-1}{2}$.
Let $n_i = |V(G_i)|$ and $e_i = |E(G_i)|$. We have $e_0 > \frac{k-1}{2}n_0$.
By induction, $e_{i+1} \geq e_i - \frac{k-1}{2} > \frac{k-1}{2}n_i - \frac{k-1}{2} = \frac{k-1}{2}(n_i - 1) = \frac{k-1}{2}n_{i+1}$.
Thus, the average degree of $G_{i+1}$ remains strictly greater than $k-1$.
Since $G$ is finite, this process must terminate at some subgraph $G^*$.
The termination condition is that every vertex in $G^*$ has degree strictly greater than $\frac{k-1}{2}$.
Since degrees are integers, $\delta(G^*) \geq \lfloor \frac{k-1}{2} \rfloor + 1 \geq \frac{k}{2}$.
Because $G^*$ is a subgraph of $G$ and $G$ does not contain $T$, $G^*$ also does not contain $T$.
Hence, $G^*$ is a minimal counterexample with $\delta(G^*) \geq \frac{k}{2}$.
This contradicts our assumption that no such subgraph exists. Therefore, such a subgraph must exist.
\end{proof}

\begin{lemma}[Path Embedding]
Let $G$ be a graph with minimum degree $\delta(G) \geq m$. Then $G$ contains a path of length at least $m$.
\end{lemma}

\begin{proof}
Let $P = (v_0, v_1, \dots, v_l)$ be a path of maximum length in $G$.
All neighbors of the vertex $v_0$ must lie on the path $P$, because if there were a neighbor $u$ of $v_0$ not in $P$, the path $(u, v_0, v_1, \dots, v_l)$ would be a longer path in $G$, contradicting the maximality of $P$.
Since $v_0$ has at least $\delta(G)$ neighbors, and all these neighbors are among $\{v_1, \dots, v_l\}$, we must have $l \geq \delta(G) \geq m$.
The path $P$ contains $l$ edges, so its length is at least $m$.
This completes the proof.
\end{proof}

\end{document}
